% ==============================================================================
% Plantilla Institucional de Trabajo Terminal — UPIIZ IPN
% Archivo: frontmatter/resumen.tex
% ==============================================================================

\chapter*{Resumen}
\addcontentsline{toc}{chapter}{Resumen}

En este proyecto se diseñará e implementará un sistema robótico móvil tipo rover autónomo capaz de explorar terrenos irregulares, desplazarse sin depender de posicionamiento satelital, identificar muestras geológicas simuladas, recolectarlas, transportarlas y depositarlas en una zona predeterminada, mientras se supervisa el estado de la misión mediante una interfaz hombre-máquina. Se desarrollará una arquitectura mecatrónica integral que articulará un sistema de locomoción todoterreno para superar desniveles y pendientes, etapas de potencia y sensado inercial y visual, algoritmos de localización relativa y evasión de obstáculos locales, así como un mecanismo automatizado de recolección y liberación de muestras sólidas. El desempeño del sistema robótico se validará experimentalmente en un entorno terrestre controlado mediante métricas cuantitativas de movilidad, precisión de navegación relativa, tasa de éxito en la recolección de muestras y porcentaje de cumplimiento integral de la misión.

\vspace{1.5em}
\noindent\textbf{Palabras clave:} Exploración autónoma, navegación sin GPS, recolección de muestras, robótica móvil, robot rover, terreno irregular.
